\section{A selector family with short covariance certificates}
\label{sec:selector}

Fix $k\ge8$, put $M=2^k$ and $d=2k+2$, and use Boolean inputs
$\xi=(x,y,z,s)\in\{-1,1\}^{2k+2}$.  The two $k$-bit blocks encode
\[
 X=\sum_{i=0}^{k-1}2^i\frac{x_i+1}{2},\qquad
 Y=\sum_{i=0}^{k-1}2^i\frac{y_i+1}{2}.
\]
For $0\le t\le M-1$, define the homogeneous affine scores
\begin{equation}\label{eq:selector-score}
 A_t(\xi)=-ts-(t+1)z+\sum_{i=0}^{k-1}2^i(x_i-y_i),
 \qquad
 Q_t=\|A_t\|_2^2=\frac{2(M^2-1)}3+2t^2+2t+1.
\end{equation}
Here $\|A_t\|_2$ denotes the Euclidean norm of its coefficient vector.
Every score is an odd integer on the full cube. On $s=1,z=-1$ it is
$1+2(X-Y)$, and on $s=z=1$ it is $2(X-Y)-2t-1$.
The target is
\begin{equation}\label{eq:selector-target}
 f(\xi)=\operatorname{sign} A_{M-1}(\xi).
\end{equation}
It is odd. On $s=1$, it is positive exactly when $z=-1$ and $X\ge Y$.

\subsection{Small pairwise-independent free-bit sets}

For $r\ge1$, let $\ell=\lceil\log_2(r+1)\rceil$ and choose $r$ distinct
nonzero vectors $v_1,\ldots,v_r\in\mathbb F_2^\ell$; the first $r$
nonzero binary vectors in their usual integer order will do. Let
\[
 \mathcal P_r=
 \left\{\bigl(-(-1)^{v_j\cdot\omega}\bigr)_{j=1}^r:
                  \omega\in\mathbb F_2^\ell\right\}.
\]
Use the uniform seed law, and let $\mathcal P_0$ consist of the single
empty assignment. The chosen vectors span $\mathbb F_2^\ell$, so the
seed map is injective. Thus
\begin{equation}\label{eq:pi-size}
 n_0:=|\mathcal P_0|=1,\qquad
 n_r:=|\mathcal P_r|=2^{\lceil\log_2(r+1)\rceil}\le2r
 \quad(r\ge1).
\end{equation}
Each free bit has mean zero, and distinct free bits have product mean
zero, because $v_i$, $v_j$, and $v_i+v_j$ are nonzero. The all-minus
assignment belongs to $\mathcal P_r$, from the zero seed. Consequently
these sets reproduce every first and second moment of independent
uniform free bits, while using only $O(r)$ points.

Define the following groups, all on $s=1$.
\begin{enumerate}
\item The positive diagonal group $\mathcal D$ has $z=-1$, $x=y$,
      with the $k$ equal bits ranging over $\mathcal P_k$.
\item For each $i=0,\ldots,k-1$, the negative carry group
      $\mathcal C_i$ has $z=-1$,
      \[
      x_j=1,\ y_j=-1\ (j<i),\qquad x_i=-1,\ y_i=1,
      \]
      and its $k-1-i$ upper equal pairs range over
      $\mathcal P_{k-1-i}$. Thus $X-Y=-1$.
\item For $k\le t\le M-1$, let $b_i=\operatorname{bit}_i(t)$.
      The negative query group $\mathcal G_t$ has $z=1$,
      \[
       (x_i,y_i)=(1,-1)\quad\text{if }b_i=1,
      \]
      and its remaining equal pairs range over
      $\mathcal P_{k-\operatorname{popcount}(t)}$. Hence $X-Y=t$.
      The point $q_t$ with $X=t,Y=0,z=s=1$ belongs to this group.
\end{enumerate}
All stated labels agree with the target. With every labeled point
$(\xi,y)$ also include its antipode $(-\xi,-y)$. Let $H_t$ be the history
containing the paired groups $\mathcal D$, all $\mathcal C_i$, and
$\mathcal G_k,\ldots,\mathcal G_t$.

\begin{lemma}[Affine minimum and covariance certificate]
\label{lem:selector-certificate}
For every $k\ge8$ and $k\le t\le M-1$, the unique minimum-norm affine
unit-margin vector fitting $H_t$ is $(0,A_t)$, where the first coordinate
is the coefficient of the appended constant. It has squared norm $Q_t$
and absolute unit margin on the whole cube. It admits an active KKT
probability distribution whose signed augmented inputs $Z$ satisfy
\[
 \mathbb E Z=(0,A_t)/Q_t,\qquad
 Z\cdot(0,A_t)=1,\qquad
 \mathbb E ZZ^{\mathsf T}\preceq2I.
\]
\end{lemma}

\begin{proof}
Write
\[
 T=t+\tfrac12,\qquad B=\frac{M^2-1}{3},\qquad
 Q=Q_t=2B+2T^2+\tfrac12,
\]
and set
\[
 c_i=2^i+Tb_i,\qquad
 S_i=\sum_{j=i}^{k-1}2^{j-i}c_j,\qquad S_k=0,
 \qquad \mu_i=c_i+S_{i+1}.
\]
These positive numbers obey
\[
 S_i=\sum_{r=i}^{k-1}\mu_r,\qquad
 S_0=B+Tt=\frac{Q-T-1/2}{2}.
\]
Give $\mathcal D$ total KKT weight $P=S_0+1/2$, each $\mathcal C_i$
total weight $\mu_i$, and $\mathcal G_t$ total weight $T$. Within each
group distribute the weight according to its uniform free-bit law.
These weights sum to $Q$. Every point in their support has signed
$A_t$-margin one.

For the moment calculations use orthonormal coordinates
\[
 d_i'=\frac{x_i-y_i}{\sqrt2},\qquad
 h_i'=\frac{x_i+y_i}{\sqrt2},\qquad
 u=\frac{s-z}{\sqrt2},\qquad v=\frac{s+z}{\sqrt2}.
\]
The coefficient vector of $A_t$ has components
\[
 (\sqrt2\,2^i)_i,\quad (0)_i,\quad 1/\sqrt2,\quad-\sqrt2 T.
\]
The signed KKT weighted mean of the points has these same components:
\[
 \sqrt2(\mu_i-S_{i+1}-Tb_i)=\sqrt2\,2^i,\qquad
 \sqrt2(P-S_0)=1/\sqrt2,\qquad -\sqrt2 T.
\]
The free-bit means vanish. This proves stationarity for the homogeneous
coefficient vector. Split every point's weight equally between it and
its labeled antipode. In the augmented coordinates this cancels the
constant-coordinate mean and leaves the other signed means unchanged.
Hence $(0,A_t)$ has a positive active KKT certificate and is the unique
minimum of the strictly convex squared norm. It fits all earlier
$\mathcal G_j$ because its score there is $2j-2t-1\le-1$.

It remains to bound the unsigned covariance $C=\mathbb E\xi\xi^{\mathsf T}$
of the normalized weights. Pairwise independence of the free equal
bits makes the $h'$ block diagonal with entries at most two; it also
makes all its cross terms with the other coordinates zero. Put
$\delta=T/Q$. The remaining block, in the order $(d',u,v)$, is
\begin{equation}\label{eq:selector-cov-block}
 C_0=\begin{pmatrix}
       K & h_u & h_v\\
       h_u^{\mathsf T}&2(1-\delta)&0\\
       h_v^{\mathsf T}&0&2\delta
     \end{pmatrix},
\end{equation}
where
\begin{align*}
 h_u&=-2c/Q,&h_v&=2\delta b,\\
 K_{ii}&=2S_i/Q+2\delta b_i,&
 K_{ij}&=-\frac2Q\bigl[2^j+Tb_j(1-b_i)\bigr]\quad(i<j).
\end{align*}
Only first and second free-bit moments enter these formulas.

Since $T<M$, $Q\ge M^2/2$, and $S_i\le2^{-i}S_0<Q/2^{i+1}$,
\[
 \delta\le2/M,\qquad \|K\|_{\mathrm{op}}\le1+6k/M.
\]
For the latter estimate, each diagonal entry is at most $1+4/M$ and
each off-diagonal absolute value is at most $6/M$; apply Gershgorin.
Furthermore, $\sum_i c_i^2\le2B+2T^2k$, $B/Q\le1/2$, and $T\ge k$
give
\[
 \frac{\|h_u\|^2}{2\delta}\le\frac2k+\frac{8k}{M},\qquad
 \frac{\|h_v\|^2}{2(1-\delta)}\le\frac{16k}{M^2}.
\]
The Schur complement of the last two coordinates in $2I-C_0$ is
\[
 2I-K-\frac{h_uh_u^{\mathsf T}}{2\delta}
       -\frac{h_vh_v^{\mathsf T}}{2(1-\delta)}.
\]
It is positive definite because
\begin{equation}\label{eq:selector-schur-bound}
 1+\frac2k+\frac{14k}{2^k}+\frac{16k}{4^k}<2
 \qquad(k\ge8).
\end{equation}
At $k=8$ the three positive terms after one sum to less than $0.69$,
and each decreases thereafter. Thus $C\preceq2I$. Pairing the augmented
inputs makes their covariance $\operatorname{diag}(1,C)$, completing
the certificate.
\end{proof}

\subsection{The fixed order and the bounded teacher}

Start the input permutation with $H_k$, in any fixed order. At stage
$t=k+1,\ldots,M-1$, present $q_t$ first, then its antipode, then the
remaining points of $\mathcal G_t$ and their antipodes in a fixed order.
At the end of the stage the history is $H_t$. Different groups are
distinguished by selector face, value of $X-Y$, or highest differing bit,
and the
antipodes have $s=-1$, so there are no repeated inputs. Append all
remaining cube points in any fixed order.

From \eqref{eq:pi-size}, the initial calibration has at most
\begin{equation}\label{eq:calibration-count}
 2\left(n_k+\sum_{r=0}^{k-1}n_r+
               n_{k-\operatorname{popcount}(k)}\right)
 \le2k^2+6k+2
\end{equation}
points. Each subsequent stage contains at most $4k$ inputs, including
its designated first query.

Before $q_t$, Lemma~\ref{lem:selector-certificate} gives the affine
minimizer $A_{t-1}$. Its score at $q_t$ is $+1$, whereas the target label
is negative. The exact profile and unique-function conclusion of
Theorem~\ref{thm:exact}, together with the posterior comparison below,
therefore give a cold-posterior mistake on every designated
first query, provided
\[
 d=2k+2\ge4\max\{1,\epsilon^{-6}\}.
\]
There are
\begin{equation}\label{eq:forced-error-count}
 M-k-1=2^k-k-1\ge2^{k-1}
\end{equation}
such queries. The unique minimizing function at the preceding history
handles all competitors, including functions that would assign the
same label to $q_t$ and its as-yet-unrevealed antipode. The posterior
argument thus does not assume that the learner's function is odd.

The target has Boolean circuit size $\Theta(k)$. A comparison circuit
computes $X\ge Y$ with $O(k)$ bounded-fan-in gates, followed by the
selector conjunction. Its odd extension on $s=-1$ is computed by
negating the original inputs conditionally on $s$ and multiplying the
output sign by $s$, again using $O(k)$ gates. Conversely, all $2k$ main
bits are essential, which gives the matching bounded-fan-in lower bound.

Finally set $C_\epsilon=\lceil1/\epsilon\rceil$. The vector $A_{M-1}/M$
has every coefficient of magnitude at most one, and has nonzero score
at every input. Use $C_\epsilon$ identical hidden neurons with this
incoming vector, zero appended-input bias, and readout weights one.
Their signed output on the target is strictly greater than
$C_\epsilon\epsilon\ge1$. This is a teacher of the same depth, width
$C_\epsilon$ independent of $k$, and individual weights at most one.
In particular the online result applies at every learner width above
this bounded teacher complexity, as well as at every other positive
width covered by the exact-profile theorem.
